%========================== METODOLOGIA ==========================
\section{Metodologia}

\subsection{Classificação da pesquisa}

Este trabalho classifica-se como pesquisa aplicada, orientada à investigação de
um problema real seguida do desenvolvimento de uma aplicação prática (pesquisa e
desenvolvimento), com o objetivo de produzir um sistema funcional de
reconhecimento de Libras a partir da investigação de técnicas de visão
computacional e aprendizado profundo já consolidadas na literatura, adaptadas ao
contexto específico da língua de sinais brasileira.

\subsection{Protocolo de coleta de dados}
\label{subsec:coleta}

O \textit{dataset} utilizado neste trabalho foi construído a partir de gravações
próprias, seguindo o protocolo descrito a seguir:

\begin{itemize}[leftmargin=2cm]
  \item \textbf{Participantes:} 10 pessoas, com no mínimo 10 repetições
  executadas por gesto;
  \item \textbf{Ambiente:} fundo neutro e iluminação adequada, minimizando ruído
  visual na etapa de detecção de \textit{landmarks};
  \item \textbf{Duração:} sessões de captura de aproximadamente 1 minuto por
  gesto por participante;
  \item \textbf{Ferramenta:} \textit{script} de coleta interativo desenvolvido
  para este trabalho (\texttt{collect\_data.py}), que grava os \textit{landmarks}
  extraídos pelo MediaPipe em formato \texttt{.npy}.
\end{itemize}

Esse protocolo é compatível, em ordem de grandeza, com \textit{datasets} de
referência da literatura de Libras: o MINDS-Libras, por exemplo, foi construído
com 12 sinalizantes cobrindo 20 sinais (REZENDE; ALMEIDA; GUIMARÃES, 2021),
reforçando que a escala adotada é adequada para um trabalho de conclusão de curso
em estágio de validação de conceito, embora ainda distante da escala necessária
para generalização robusta (ver Seção~\ref{subsec:limitacoes}).

\begin{figure}[htbp]
\centering
\figpendente{Print da tela de coleta: \textit{landmarks} desenhados sobre a mão,
gesto/participante atual e contador de repetições.}
\caption{Interface de coleta de dados (\texttt{collect\_data.py}) durante a
gravação de um gesto.}
\label{fig:coleta}
\end{figure}

\subsection{Critérios de exclusão e rotulação das amostras}

Foram descartadas do \textit{dataset} final as amostras em que os
\textit{landmarks} da mão eram perdidos durante a captura — por exemplo, quando a
mão saía parcialmente do quadro ou o rastreamento do MediaPipe falhava
momentaneamente —, bem como gravações em que o gesto foi executado de forma
incompleta. Esse critério de exclusão foi aplicado automaticamente durante a
etapa de coleta, descartando a sequência em tempo real sempre que a detecção de
\textit{landmarks} era interrompida.

% [PENDENTE: descrever como as amostras foram rotuladas — manualmente ou
% automaticamente a partir da estrutura de pastas do dataset (dataset/<gesto>/...)]

\subsection{Pipeline de processamento}
\label{subsec:pipeline}

O sistema opera em um \textit{pipeline} sequencial de sete etapas, executado
inteiramente de forma local, sem servidor ou \textit{backend} remoto. A
Figura~\ref{fig:pipeline} apresenta o fluxo, e a Tabela~\ref{tab:pipeline}
detalha cada etapa.

\begin{figure}[htbp]
\centering
\figpendente{Diagrama do \textit{pipeline}: 1.~Captura (OpenCV, 640$\times$480) →
2.~Detecção (MediaPipe, 21 \textit{landmarks}) → 3.~Normalização (centro no
pulso) → 4.~\textit{Buffer} (50 \textit{frames}, passo 10) → 5.~CNN + LSTM
(entrada 50$\times$63) → 6.~Limiar (confiança $\geq$ 70\%) → 7.~Frase
(\textit{debounce} 1,5~s; WORDS / LETTERS / MIXED).}
\caption{Pipeline de processamento do LibrasAI, da captura à formação de frases.}
\label{fig:pipeline}
\end{figure}

\begin{table}[htbp]
\centering
\caption{Etapas do pipeline de processamento.}
\label{tab:pipeline}
\small
\begin{tabular}{@{}c p{3cm} p{7.5cm}@{}}
\toprule
\textbf{\#} & \textbf{Etapa} & \textbf{Detalhamento técnico} \\
\midrule
1 & Captura &
OpenCV abre a webcam com \textit{backend} \texttt{CAP\_DSHOW} (Windows).
Aplica-se espelhamento horizontal e redimensionamento para 640$\times$480. \\
\addlinespace
2 & Detecção de mãos &
MediaPipe Hands processa o \textit{frame} RGB e retorna 21 \textit{landmarks}
$(x, y, z)$ normalizados por mão, com suporte a 2 mãos simultâneas. \\
\addlinespace
3 & Normalização espacial &
Centralização dos \textit{landmarks} pelo pulso (\textit{landmark} 0) e divisão
pela distância pulso–base do dedo médio (\textit{landmark} 9), removendo variação
de posição e de distância da câmera. \\
\addlinespace
4 & \textit{Buffer} deslizante &
\textit{Buffer} circular de 50 \textit{frames} (\texttt{SequenceBuffer}). A cada
10 novos \textit{frames} (\texttt{step\_size}), uma inferência é disparada,
formando uma janela deslizante para predição contínua. \\
\addlinespace
5 & Modelo CNN + LSTM &
Entrada $(1, 50, 63)$. \texttt{TimeDistributed(Conv1D)} extrai características
espaciais por \textit{frame}; \texttt{LSTM(128)} e \texttt{LSTM(64)} modelam a
dinâmica temporal; \texttt{Dense(64)} e \textit{softmax} classificam. \\
\addlinespace
6 & Limiar de confiança &
Apenas predições com confiança $\geq$ 70\% são aceitas. As três classes mais
prováveis são exibidas como barras de probabilidade na tela. \\
\addlinespace
7 & Formação de frases &
\textit{Debounce} temporal de 1,5~s evita repetições indevidas. Modos: WORDS
(gesto = palavra), LETTERS (gesto = letra, acumulado em palavra) e MIXED (seleção
automática por tamanho). \\
\bottomrule
\end{tabular}
\end{table}

\subsection{Arquitetura do modelo}
\label{subsec:arquitetura}

A rede neural híbrida recebe como entrada um tensor de dimensões $(50, 63)$ — 50
\textit{frames} por 21 \textit{landmarks} por 3 coordenadas. Camadas
convolucionais distribuídas no tempo extraem características espaciais de cada
\textit{frame}; em seguida, duas camadas LSTM modelam a dinâmica temporal da
sequência, e camadas densas realizam a classificação. A
Tabela~\ref{tab:topologia} descreve a topologia camada a camada.

O treinamento foi conduzido com os \textit{callbacks} \texttt{EarlyStopping},
\texttt{ReduceLROnPlateau}, \texttt{ModelCheckpoint} e \texttt{CSVLogger}, e o
modelo final foi exportado tanto em formato TFLite, para viabilizar uso futuro em
dispositivos móveis, quanto em \texttt{SavedModel}, para uso em produção.

% [PENDENTE: informar os hiperparâmetros de treinamento: otimizador, taxa de
% aprendizado inicial, batch size, número máximo de épocas e paciência do
% EarlyStopping; a divisão treino/validação/teste (proporções e se a divisão foi
% por amostra ou por participante); e as técnicas de data augmentation aplicadas.]

\begin{table}[htbp]
\centering
\caption{Topologia do modelo CNN + LSTM ($N$ = número de classes).}
\label{tab:topologia}
\small
\begin{tabular}{@{}c p{3.6cm} p{5.6cm} c@{}}
\toprule
\textbf{\#} & \textbf{Camada} & \textbf{Configuração} & \textbf{Saída} \\
\midrule
1 & Entrada & 50 \textit{frames} $\times$ 63 valores & $(50, 63)$ \\
2 & \texttt{Reshape} & — & $(50, 63, 1)$ \\
3 & \texttt{TimeDistributed(Conv1D)} & 32 filtros, \textit{kernel} 3, ReLU, \textit{padding} \texttt{same} & $(50, 63, 32)$ \\
4 & \texttt{TimeDistributed(Conv1D)} & 64 filtros, \textit{kernel} 3, ReLU, \textit{padding} \texttt{same} & $(50, 63, 64)$ \\
5 & \texttt{TimeDistributed(GlobalAvgPool1D)} & — & $(50, 64)$ \\
6 & \texttt{TimeDistributed(Dropout)} & taxa 0,2 & $(50, 64)$ \\
7 & \texttt{LSTM} & 128 unidades, \texttt{return\_sequences=True}; \texttt{BatchNormalization}; \textit{Dropout} 0,4 & $(50, 128)$ \\
8 & \texttt{LSTM} & 64 unidades, \texttt{return\_sequences=False}; \texttt{BatchNormalization}; \textit{Dropout} 0,4 & $(64)$ \\
9 & \texttt{Dense} & 64 unidades, ReLU; \textit{Dropout} 0,2 & $(64)$ \\
10 & \texttt{Dense} (saída) & $N$ unidades, \textit{softmax} & $(N)$ \\
\bottomrule
\end{tabular}
\end{table}

\subsection{Protocolo de validação com usuários}
\label{subsec:validacao}

Uma vez que o modelo foi treinado e testado apenas com dados de participantes
ouvintes envolvidos no desenvolvimento, é necessária uma etapa de validação com o
público-alvo do sistema. Uma tentativa inicial de coleta por meio de formulários
\textit{on-line} não produziu resultados práticos, sobretudo pela ausência de
mediação em Libras. Por isso, optou-se por uma validação assistida, presencial ou
remota, com o apoio de pessoas fluentes em Libras.

\subsubsection{Participantes}

A validação será conduzida com um grupo de quatro a cinco integrantes da
comunidade surda e usuária de Libras, contatados por meio de um grupo localizado
no município de Gravatá (PE). Serão incluídos participantes maiores de idade,
usuários de Libras, com qualquer grau de surdez, incluindo usuários de implante
coclear. Nenhum dos participantes terá contribuído com amostras para o conjunto
de treinamento, de modo que a validação também cumpra o papel de avaliação
independente do sinalizante (\textit{signer-independent}).

\subsubsection{Procedimento}

Cada sessão seguirá as etapas abaixo:

\begin{enumerate}[leftmargin=2cm]
  \item apresentação do estudo em Libras e assinatura do Termo de Consentimento
  Livre e Esclarecido (TCLE);
  \item registro dos dados demográficos do participante na plataforma, associados
  a um identificador anônimo (P01, P02, \dots);
  \item breve demonstração do sistema e período de familiarização;
  \item execução, pelo participante, dos gestos do vocabulário reconhecido, com
  repetições, enquanto a plataforma registra automaticamente as métricas de
  interação;
  \item formação livre de frases curtas com o vocabulário disponível;
  \item aplicação do questionário de percepção, com mediação em Libras.
\end{enumerate}

\subsubsection{Dados coletados}

A plataforma deverá registrar e associar automaticamente a cada identificador
anônimo os seguintes dados:

\begin{itemize}[leftmargin=2cm]
  \item \textbf{Demográficos:} sexo, faixa etária, grau de surdez, uso de
  implante coclear ou aparelho auditivo e tempo de uso de Libras;
  \item \textbf{Métricas de interação:} gesto esperado e gesto predito, confiança
  da predição, tempo até o reconhecimento, número de tentativas por gesto e
  frases formadas;
  \item \textbf{Percepção:} respostas ao questionário de \textit{feedback},
  contemplando facilidade de uso, adequação dos gestos reconhecidos, velocidade
  de resposta e utilidade percebida, em escala visual (por exemplo, de 1 a 5) e
  com espaço para comentários registrados com auxílio do intérprete.
\end{itemize}

Como parte do público-alvo pode ter proficiência limitada em português escrito, o
questionário será apresentado com apoio de intérprete e com escalas visuais, em
vez de perguntas exclusivamente textuais.

\subsubsection{Análise}

A partir das métricas registradas serão calculados, por participante e no
agregado, a acurácia de reconhecimento, a matriz de confusão e o tempo médio até
o reconhecimento. As respostas do questionário serão sumarizadas por estatística
descritiva. Os resultados serão comparados aos obtidos no conjunto de teste
original (Seção~\ref{subsec:metricas-modelo}), permitindo estimar a perda de
desempenho em usuários externos ao treinamento.

\subsubsection{Aspectos éticos}

Os dados serão armazenados de forma anonimizada, sem gravação de vídeo
identificável: a plataforma armazena apenas as coordenadas dos \textit{landmarks}
e as métricas derivadas.

% [PENDENTE: confirmar com a coordenação se o estudo precisa ser submetido a
% Comitê de Ética em Pesquisa (Plataforma Brasil) e anexar o TCLE como apêndice.]
